% Task-18 full methods, moved from the main text. Summary in lewm18-methods.tex.
\section{PushT and Cube Methods}\label{app:extension-methods}
\subsection{Tasks, cases and methods}
In Cube, success means the cube's centre comes within 4 cm of the target.
(An earlier Cube probe in our records tested only gripper reaching, so it is
not evidence for this goal.) Each task uses its own frozen LeWM checkpoint.
Each case takes a start image and a goal image 25 recorded actions apart in
the task's pinned demonstration dataset, places a forbidden band across the
straight line from start to goal with an off-centre gap, and keeps the start
and goal outside the enlarged danger region. Development and
final cases come from different demonstrations. A 0 margin keeps only the
original danger region. Native planning ignores the margin, so its attempts
are recorded once and scored at all three margins; these are three views of
the same attempts, not three independent samples.

The allowance is 50 executed actions, with a replan after every five.
Original SafeCEM uses 30 CEM iterations, 300 candidates and five blocks of
look-ahead. Shorter look-ahead receives more candidates under the same
allowance of 1,500 predictor steps per iteration. We report actual model
work separately, because early completion, failed replans and the extra
prediction for the chosen plan change the work used.

\subsection{Selection}
The Feedback controller uses latent progress and prediction error to change look-ahead, sampling
spread and candidate reuse; it cannot weaken the margin. For each margin, six
fixed configurations are compared on six development cases with two seeds.
Four Feedback presets then use the same six cases and one seed. Selection
first maximizes safe completion, then breaks ties by fewer failure-capped
actions, less model work and finally list order. Feedback builds on the
fixed-setting selection and uses another 72 development attempts per task,
so its tuning cost is higher, not equal.

\subsection{What the checks and speed measurements mean}
The safety rule covers only the pusher centre in PushT and the gripper pinch
point in Cube. It does not certify the moved object, the arm or the whole
robot against collision. In PushT, the checker uses the known pusher
dynamics and tests every physics step. In Cube, it ranks candidates with an
approximate path, then tests the first five actions of up to eight
candidates in the true simulator. Every such preview counts toward the work
totals, and the full simulator state is restored before the next preview or
the real action. If no preview passes, the attempt ends as a controller
failure; there is no unchecked fallback.

PushT keeps the pinned environment's own success test, which requires the
pusher position, block position and block angle to be close to the goal; we
do not report shape overlap. Cube's goals come from LeWM's dataset pairs, not
the full official five-goal OGBench evaluation. As a small protocol check,
native planning also ran on twelve separate pairs per task executing all
25 planned actions before replanning, as the reference planner does, without
a danger region or checker. It completes 11/12 PushT pairs and 10/12 Cube
pairs, but seven of the Cube pairs already meet the object goal at reset, so
this panel must not replace the reserved cases.

Speed is scored in failure-capped actions (Section~\ref{sec:success}): a safe
completion costs the actions it used, and every unsafe or unfinished attempt
costs 50, so early failure cannot look fast. Mean actions among successful
attempts are secondary and are reported with their own counts. A separate
run, one attempt at a time, measures actual planning time on the first two
final cases at the largest margin, with a new shared seed. Video playback
speed says nothing about computation time. Uncertainty intervals resample
whole cases, keeping each case's seeds and paired planners together.

\subsection{The PushT reset repair}
A review before publication found that resetting PushT moves the simulator
slightly, but the checker in our first cohort started from the state before
the reset. We made the checker start from the state that reset returns and
repeated the full protocol, including development selection, with the same
cases and candidate configurations. The first cohort's 676 attempts are kept
and counted in the work totals, but excluded from the performance results.
